%==============================================================================
% Sjabloon onderzoeksvoorstel bachproef
%==============================================================================
% Gebaseerd op document class `hogent-article'
% zie <https://github.com/HoGentTIN/latex-hogent-article>

% Voor een voorstel in het Engels: voeg de documentclass-optie [english] toe.
% Let op: kan enkel na toestemming van de bachelorproefcoördinator!
\documentclass{hogent-article}
\setlength{\emergencystretch}{3em}

% Invoegen bibliografiebestand
\addbibresource{SaricMilan-bpvoorstel.bib}

% Informatie over de opleiding, het vak en soort opdracht
\studyprogramme{Professionele bachelor toegepaste informatica}
\course{Bachelorproef}
\assignmenttype{Onderzoeksvoorstel}


\academicyear{2025-2026} % TODO: pas het academiejaar aan

% TODO: Werktitel
\title{Proof of concept: model context protocol onderstuende coding agent voor game server management.}

% TODO: Studentnaam en emailadres invullen
\author{Milan Saric}
\email{milan.saric@student.hogent.be}


% TODO: Geef de co-promotor op
\supervisor[Co-promotor]{N. Candaele (Takaro.io, \href{mailto:niek.candaele@gmail.com}{niek.candaele@gmail.com})}

% Binnen welke specialisatierichting uit 3TI situeert dit onderzoek zich?
% Kies uit deze lijst:
%
% - Mobile \& Enterprise development
% - AI \& Data Engineering
% - Functional \& Business Analysis
% - System \& Network Administrator
% - Mainframe Expert
% - Als het onderzoek niet past binnen een van deze domeinen specifieer je deze
%   zelf
%
\specialisation{AI \& Data Engineering}
\keywords{LLM, MCP, POC}

\begin{document}

\begin{abstract}
Large Language Models (LLM’s) demonstreren indrukwekkende capaciteiten in het genereren van code, maar het correct aanroepen en toepassen van complexe, domeinspecifieke API’s blijft een uitdaging. Het Model Context Protocol (MCP) biedt een gestructureerde manier om API-functionaliteiten op een gecontroleerde wijze beschikbaar te maken voor LLM’s. Game server managementplatformen, zoals Takaro.io, beschikken over uitgebreide API-endpoints, wat het genereren van correcte en herbruikbare codefragmenten bemoeilijkt door de beperkte context die overblijft na het toevoegen van MCP-servers. Dit onderzoek heeft als doel het ontwikkelen van een proof of concept (POC) van een MCP-ondersteunde codegeneratie-agent die, op basis van gebruikersprompts, correcte codefragmenten genereert voor game serverbeheer binnen een specifiek contextvenster. Na validatie door de gebruiker kan de gegenereerde code worden toegevoegd als herbruikbare module binnen het platform. Tot slot evalueert dit onderzoek in welke mate MCP bijdraagt aan zowel de correctheid van de gegenereerde code als de resterende beschikbare context.
\end{abstract}

\tableofcontents

% De hoofdtekst van het voorstel zit in een apart bestand, zodat het makkelijk
% kan opgenomen worden in de bijlagen van de bachelorproef zelf.
\input{voorstel-inhoud}

\printbibliography[heading=bibintoc]

\end{document}